\section{Cyclotomic coordinates, evaluation, and integral Herr classes}
\label{sec:wach-v059}

Fix an original prime $p$, local period datum $P$, and good layer
$(K,s,L)$. Write $O=\mathcal O_L$, $K_0$ for the maximal unramified
subfield of $K$, $n=[K_0:\mathbb Q_p]$, $J=\mathbb Z/n\mathbb Z$,
and $G=G_K$. The index $i-1$ below belongs to $J$. We use $C$ for
the completed algebraic closure and $C^\flat$ for its tilt. The group
record determines which Frobenius and Galois substitutions act on the
series ring. The evaluation record in turn identifies the coefficient
rings on which FET25 constructs the original \'etale functor. We then
construct the two integral class maps of that functor.

% BEGIN LWR1
\paragraph{LWR1. The actual group and its two distinguished elements.}
\label{claim:LWR1}
Every good layer admits the original $\mathrm{GamS}$, including the
tie rule for the two generators and the entire finite torsion subgroup.

\begin{proof}
Use WPR2 and put $H=\chi(G)\subset\mathbb Z_p^\times$. This is compact
and contains $U_s=1+p^s\mathbb Z_p$ by goodness and WPR2. Set
$\Gamma=H$, use the range restriction of $\chi$ for the quotient map,
and embed $H$ into $O^\times$ by the coefficient inclusion. This is
injective and its composite is the supplied period character. The
original valuation is normalized by $v(p)=1$ and has valuation ring
$O$. Its comparisons on $L$ therefore agree with the usual
$p$-adic ones: to test $v(x)\ge a/b$ with $b>0$, test integrality
of $x^bp^{-a}$. Thus $G_s=\chi^{-1}(U_s)$ has the required defining
inequality. Goodness gives $s\ge1$, and $s\ge2$ for $p=2$.

Let $\gamma=1+p^s$ and $\Delta=H_{\mathrm{tors}}$. The logarithm
and exponential are inverse isometries between $U_1$ and $p\mathbb Z_p$
for odd $p$, and between $U_2$ and $4\mathbb Z_2$ at 2. This follows
from their convergent series: the linear term strictly dominates the
remaining terms on these balls, and the formal inverse identities
converge there. Thus reduction modulo $p^2$ injects $\Delta$ into a
finite group. Set $d=|\Delta|$. Among all $h\in H$ minimize
$v(h^d-1)$, choosing $\gamma_K=\gamma$ if it is a minimizer.
The minimum exists by discreteness and is finite since $\gamma$ has
infinite order. This gives the precise valuation inequality and tie
rule. For $s=1$ and odd $p$, the Teichm\"uller lift of the reduction
of every $h\in H$ is in $H$ since its quotient by $h$ is in $U_1$.
Reduction injects $\Delta$ into $\mathbb F_p^\times$, so $p\nmid d$,
all $h^d$ are $1$ modulo $p$, and
$(1+p)^d\equiv1+dp\pmod{p^2}$. Hence $\gamma$ is a minimizer here.

Let $\iota$ be the inclusion of $K_0$ in $L$, available by goodness.
Write $\sigma$ for the specified Frobenius of the period datum. Its
fixed field on $K_0$ is $\mathbb Q_p$: a fixed element is Frobenius
fixed in the crystalline period ring and belongs to filtration zero
in the de Rham ring; the fundamental period identity in $P$ puts it
in $\mathbb Q_p$. The period-ring embedding of the field is injective
(the period parameter and the filtration axiom give a nonzero ring).
The finite group generated by $\sigma$ consequently has order $n$ by
the fixed-field theorem. The $n$ distinct embeddings
$e_i=\iota\sigma^{-i}$ exhaust the embeddings in $L$ and
$e_{i+1}\sigma=e_i$. Finally use the line $L$ with trivial action
and lattice $O\cdot1$. These choices give all fields of the record.
\end{proof}
% END LWR1

% BEGIN LWR2
\paragraph{LWR2. The original series ring and every level evaluation.}
\label{claim:LWR2}
For every supplied $\Gamma$ on this layer there is a $\mathrm{WSer}(\Gamma)$.
In particular its evaluation law holds coefficient by coefficient at
every precision, including evaluation away from primitive roots.

\begin{proof}
For any such $\Gamma$ its quotient character has a unique lift
$c:\Gamma\to\mathbb Z_p^\times$. Lift an element to $G$, apply WPR2,
and use the character identity in $L$ and injectivity of
$\mathbb Z_p\hookrightarrow L$ to see independence and multiplicativity.
Put $A=\prod_{i\in J}O[[X]]$ and $\pi=(X)_i$. On each copy take
\[
 X\longmapsto(1+X)^p-1=Xq(X),\quad q(X)=\sum_{a=0}^{p-1}(1+X)^a,
 \qquad X\longmapsto (1+X)^{c(g)}-1.
\]
WPR3 gives composition of binomial substitutions, their commutation,
and the specified linear and constant terms. Frobenius shifts the
input letter to $i-1$; the Galois substitution preserves the letter.
Use the $p$-adic completion of $A[1/\pi]$ for $E_E$. Modulo $p^k$
its elements have Laurent representatives with one common pole bound
since $J$ is finite. The map from $A/p^k$ is injective, as $X$ is a
nonzerodivisor even over $O/p^k$. It follows that $A\to E_E$ is
injective, its image is $p^k$-saturated, and for every element of
$E_E$ a power of $\pi$ clears the poles modulo $p^k$. Here the
inverse-limit description gives $E_E/p^k=A[1/\pi]/p^k$: the
localized ring has no $p$-torsion, so compatible classes zero modulo
$p^k$ can be divided by $p^k$ in the inverse system. The two operators
extend: the image of $X$ under the first is $X^p$ modulo $p$ and
under the second is $X$ times a series with unit constant term.
Both images become units. Uniqueness on the localization and density
give all extension and composition laws.

If $\zeta\in L$ has $p$-power order, set
$x=\zeta\exp(t)-1$ and $b=\zeta-1$. Reduction in the residue field
gives $\bar\zeta=1$, so $|b|<1$. Define
\[
    {\rm lev}_{i,\zeta}(f)=f_i(\zeta\exp(t)-1)\in L[[t]].
\]
For arbitrary $x$ with constant term $b$, first translate $X$ to
$X+b$ in $O[[X]]$, then substitute $x-b\in tL[[t]]$. Translation
is well defined since $b$ lies in the maximal ideal:
the $j$th coefficient of the translated series
$\sum a_m X^m$ is $\sum_{m\ge j}\binom mj a_m b^{m-j}$.
It is an algebra automorphism, with inverse translation by $-b$
(equivalently use $(\mathfrak m_O,X)$-adic completeness). This also
proves uniqueness without any continuity hypothesis on a competing
algebra map: division by $X-b$ with constant remainder can be
iterated, and agreement modulo every power of $t$ forces equality.
For fixed $k$, replacing $f_i$ by its truncation of degree $<m$
changes each of the first $k+1$ translated coefficients by at most
$|b|^{m-k}$ when $m>k$ (at $b=0$ these changes vanish).
Substitution of $x-b$ multiplies these bounds by a finite constant
depending on $x,k$. This gives the required valuation threshold for
every $r$. For other $\zeta$ use the evaluation at $1$.

If $\zeta$ has order $p^j$ and $c\equiv a\pmod{p^j}$, then
\[
 (1+X)^c\big|_{X=\zeta e^t-1}=\zeta^a e^{ct}.
\]
Indeed choose nonnegative integers $c_N\to c$ with the same eventual
residue modulo $p^j$. Binomial coefficients of each fixed order
converge in $\mathbb Z_p$. The preceding tail bound is uniform over
integral series, so each evaluated coefficient converges as well.
The identity for integer powers passes to the limit.
Uniqueness of evaluation now extends the identity on $X$ to all
series. Thus $\mathrm{lev}_{i,\zeta}(g f)$ is
$\mathrm{lev}_{i,\zeta^a}(f)$ with $t$ replaced by $c(g)t$.
The original inequality $j\le v(c(g)-a)$ is precisely the required
congruence. Frobenius gives $\mathrm{lev}_{i-1,\zeta^p}(f)$ with
$t$ replaced by $pt$. The image of $1+\pi$ is $\zeta e^t$, where
$(e^t)'=e^t$ and $e^0=1$. Evaluation kills a zero input letter.
These verify the remaining fields.
\end{proof}
% END LWR2

% BEGIN LWR3
\paragraph{LWR3. The actual integral inclusion and its extension.}
\label{claim:LWR3}
For an arbitrary supplied $\Gamma,R$ and local $A_{\inf}$ datum
$AI$, there are $I:\mathrm{IoS}(AI,R)$ and $\mathrm{IoW}(I)$.
This constructs the input that WPR14 and FET25 previously used.

\begin{proof}
Write $B=O\otimes_{\mathbb Z_p}W(\mathcal O_{C^\flat})$ and
$D=O\otimes_{\mathbb Z_p}W(C^\flat)$. We first specify the map on
constant tuples without using any existing inclusion $I$.
Let $W_0=\mathcal O_{K_0}$. The letter map
$O\otimes W_0\to O^J$, $a\otimes w\mapsto(a e_i(w))_i$,
is a ring isomorphism by WPR26. The residue field $k_0$ of $W_0$
has its prescribed image in the residue field of $C$ via $K_0\subset C$.
Identify this residue field with that of the tilt by its residue map.
Then $k_0$ embeds in $\mathcal O_{C^\flat}$: in characteristic $p$ the
finite field of its cardinality lifts canonically through residue
reduction, using the simple roots of $Z^{|k_0|}-Z$. Identify
$W_0=W(k_0)$ by the unramified integer-ring isomorphism
(the Teichm\"uller representatives and their convergent $p$-adic digit
expansions give this identification). Witt functoriality and the
inverse of the letter map give an $O$-algebra map $u:O^J\to B$.
For each prime-to-$p$ root $w\in K_0$ its letter tuple maps to
$1\otimes[w^\flat]$. Indeed the compatible inverse $p$-powers of $w$
in its finite group are unique; their reductions form the given
tilt, and the unique Teichm\"uller lift is the original root in $W_0$.
This gives the entire prescribed constant-tuple law.

Put $e=AI.\epsilon_s$, $y=1\otimes([e]-1)$, and
$E_i=u(\mathbf 1_i)$. The orthogonal idempotents $E_i$ sum to one.
Define $\mathrm{io}(f)=\sum_i E_i f_i(y)$. To see convergence directly,
$e-1$ is topologically nilpotent by WIS3. Write
$[e]-1=[e-1]+pb$ in the Witt ring. Modulo $p^k$, its $N$th power is
divisible by $[e-1]^{N-k+1}$ for $N\ge k$. In the first $k$ Witt
coordinates multiplication by $[e-1]^m$ multiplies coordinate $j$ by
$(e-1)^{mp^j}$. Thus uniformly for integral coefficients the tails
converge to zero in each truncated Witt ring with its product
valuation topology. Completeness follows coordinatewise from
completeness of $\mathcal O_{C^\flat}$ (addition, negation and
multiplication are continuous triangular Witt polynomials).
A finite $\mathbb Z_p$-basis of $O$ gives the same assertion for $B$.
Products pass to the limit, giving an algebra map with $\pi\mapsto y$.

On the constant tuples Frobenius takes the tuple of $w$ to that of
$\sigma w$, hence to $1\otimes[(w^\flat)^p]$ on the generating
prime-to-$p$ roots. Indeed the residue-Frobenius axiom of the supplied
period datum (also used in WPR1) gives $\sigma(w)\equiv w^p$ in the
residue field. They are roots of the same $Z^m-1$, with $m$ prime
to $p$, so uniqueness of the simple lift gives $\sigma(w)=w^p$.
The action of $G$ fixes these roots and their
compatible tilts. WPR1 supplies their integral span. On $\pi$,
Frobenius sends $[e]-1$ to $[e]^p-1$. For $g\in G$ the specified
series operator is the binomial substitution by $c=\chi(g)$ by
WPR5--WPR6, valid for every supplied $R$. The value of $(1+X)^c$
at $[e]-1$ is $[ge]$. To see this, approximate $c$ by nonnegative
integers $a_N\to c$. Tails are uniformly small in each finite set of
Witt coordinates by the preceding argument, and each fixed binomial
coefficient converges $p$-adically. The integer evaluations are
$[e^{a_N}]$. For the $j$th inverse-tilt coordinate the cyclotomic
root has order dividing $p^{s+j}$, and $a_N$ eventually agrees with
the exponent of $g$ modulo $p^{s+j}$. Agreement of the first $j$
tilt coordinates implies convergence in the tilt valuation topology
(the difference in coordinate zero after taking compatible roots
has norm at most $|p|^{p^j}$). This proves the claimed value.
The two composites therefore agree on constants and $\pi$. They
agree on all series: after truncating at degree $k$ their difference
is divisible by the $k$th power of the common image of $\pi$.
WPR13 separates these powers for both transformed parameters
$e^p$ and $ge$, whose differences from $1$ are nonunits.

In $D$ the image of $y$ is a unit since $e-1\ne0$ in $C^\flat$:
$e^{p^s}=\epsilon\ne1$. Also $D$ is $p$-adically complete,
by the finite basis of $O$. For arbitrary $x\in R.E_E$ choose
$\pi^m x\equiv a\pmod{p^k}$ from its density field and assign the
residue $y^{-m}\mathrm{io}(a)\in D/p^k$. Two choices agree:
cross-multiply and use $R.E_E$'s saturation field to see that
$a\pi^{m'}-a'\pi^m$ already belongs to $p^k A$. These residues are
compatible. The resulting $\mathrm{io}_E$ is an algebra map by
addition and multiplication of the representatives. Density,
invertibility of $y$, and separation prove uniqueness, also for the
two equivariance diagrams (the transformed images of $\pi$ are
units). This proves every field of $I$ and the indicated $\mathrm{IoW}$
law for the same map. Consequently FET25 supplies $E$ on this layer;
WPR30 supplies $\mathrm{NS}(E)$ under its original BS and
$\mathrm{LocTied}$ hypotheses.
\end{proof}
% END LWR3

% BEGIN LWR4
\paragraph{LWR4. Two integral coordinates on the actual Galois groups.}
\label{claim:LWR4}
For any supplied group record, $G_s$ has a continuous additive
$p$-adic coordinate $\lambda_s$ whose fibers at $0$ and $1$ are
exactly the fibers of the quotient map at $1$ and $\gamma$.
On $G$ there is a continuous additive coordinate $\lambda_K$ whose
kernel is the preimage of $\Delta$ and which takes value $1$ on the
fiber of $\gamma_K$. Both coordinates are onto.

\begin{proof}
As in LWR2 identify the group with $H=\chi(G)$ by its injective
coefficient character. The subgroup $G_s$ is the inverse image of
$U_s$, and $U_s\subset H$ by WPR2. On it take
$\lambda_s(g)=\log(\chi(g))/\log(\gamma)$, which is integral by
the isometry on $U_s$ and $v(\gamma-1)=s$. It takes the two
specified fibers because logarithm is injective there. Surjectivity
also follows from compact image and the value $1$.

Write $L_p$ for logarithm on the whole unit group, extended by zero
on torsion: decompose into the residue Teichm\"uller factor and
$U_1$ for odd $p$, and into $\{\pm1\}$ and $U_2$ for $p=2$.
Its kernel on $H$ is precisely $\Delta$; its image is a nonzero
compact additive subgroup of $\mathbb Q_p$ containing $p^s\mathbb Z_p$.
We claim $L_p(\gamma_K)$ generates this image over $\mathbb Z_p$,
for the supplied choice of minimizer. When $p$ is odd the torsion
factor of each $h\in H$ lies in $H$: choose positive exponents
tending $p$-adically to zero and congruent to $1$ modulo $p-1$,
and take the limit of these powers in the closed group $H$.
So all $h^d$ belong to $U_1$ for $d=|\Delta|$. At $p=2$ with
$-1\in H$ all $h^d$ belong to $U_2$. In either case the logarithm
isometry gives $v(h^d-1)=v(d)+v(L_p(h))$, with the usual infinite
value at zero. Minimization proves the assertion.

In the remaining case $p=2$ and $-1\notin H$, $d=1$ and $L_2$ is
a topological isomorphism of $H$ onto $2^a\mathbb Z_2$ for some
$a$. If $H\subset U_2$, the same isometry applies. Otherwise the
sign character $H\to\{\pm1\}$ (reduction modulo $4$) is onto.
The elements congruent to $-1$ modulo $4$ are exactly the minimizers,
since their differences from $1$ have valuation one. Their logarithms
lie outside the unique subgroup of index two in $2^a\mathbb Z_2$,
so each is a generator. This covers the exceptional dyadic case
without projecting away a sign inside the supplied $H$. Now put
$\lambda_K(g)=L_p(\chi(g))/L_p(\gamma_K)$. The resulting integral
homomorphism has the asserted fibers and is onto.
\end{proof}
% END LWR4

% BEGIN LWR5
\paragraph{LWR5. The actual Herr class maps, including integral restriction.}
\label{claim:LWR5}
For any supplied $\Gamma,R,I,\mathrm{IoW}(I),E$ on a good layer,
the original $\mathrm{WCl}(E)$ exists. Its two maps take values in
the cohomology of the given integral lattice.

\begin{proof}
Write $T$ for the underlying stable $O$-lattice and
$M=D\otimes_O T$, where $D=O\otimes_{\mathbb Z_p}W(C^\flat)$.
Let $j(t)=1\otimes t$, let $a_g$ denote the diagonal Galois action,
and put $\delta b(g)=a_g b-b$. Frobenius $\varphi$ acts on the
Witt factor, so it commutes with $a_g$. On the actual lattice
\[
  0\longrightarrow T\xrightarrow{j}M
     \xrightarrow{\varphi-1}M\longrightarrow0
\]
is exact. Indeed $C^\flat$ is algebraically closed, so $u^p-u=v$
is solvable for every $v$, with kernel $\mathbb F_p$. Correct a
Witt lift successively modulo $p,p^2,\ldots$ using this equation
on each residual digit. The limit solves the Witt equation, and a
fixed Witt vector has all coordinates in $\mathbb F_p$, giving
kernel $W(\mathbb F_p)=\mathbb Z_p$. A finite $\mathbb Z_p$-basis
of $T$ transfers exactness to $M=W(C^\flat)\otimes_{\mathbb Z_p}T$.
Such a basis exists since $T$ is a finite torsion-free lattice over
the DVR $O$, and $O$ is finite free over $\mathbb Z_p$.

Give the Witt ring the product of the valuation topologies on its
coordinates, and $M$ the resulting finite product topology in this
basis. It is Hausdorff, with continuous ring operations on the Witt
factor, since each output coordinate uses finitely many polynomial
operations. The action on $C^\flat$ is continuous: on integral
tilts each prescribed residue coordinate is locally constant in the
Galois argument (approximate its representative modulo $p$ by an
algebraic element fixed by an open subgroup); the coordinate bound
used in LWR3 gives continuity of each orbit. The action preserves
the norm, so joint continuity follows by the triangle inequality.
Extend to the fraction field, and then apply this coordinatewise
to Witt vectors. The matrix for the action on $T$ is continuous
because its basis coordinate maps extend to continuous rational
linear functionals on the original finite dimensional space.
The map $\mathbb Z_p\to W(C^\flat)$ is continuous since congruence
modulo $p^k$ gives agreement of the first $k$ coordinates.
Together these statements give continuity of $(g,m)\mapsto a_g m$.
The subgroup $j(T)$ is compact and $j$ is an embedding by compact
to Hausdorff injectivity. Also $p^k j(T)$ tends uniformly to zero.

In the restricted case use $G'=G_s$, $\lambda=\lambda_s$,
$\eta=\gamma$; in the full case use $G'=G$, $\lambda=\lambda_K$,
$\eta=\gamma_K$. The latter domain consists of Herr pairs fixed
in both coordinates by every member of $\Delta$. Put
$y=j_D(u)$ and $z=j_D(v)$ for a Herr pair $(u,v)$, using the
given embedding of $E$. Pick $g_0$ lifting $\eta$. Both $y$ and $z$
are fixed by $N=\ker\lambda$: in the first case its quotient image
is trivial and in the second it is torsion. Equivariance of $j_D$
and the Herr-pair equation give
\[
           a_{g_0}y-y=(\varphi-1)z.
\]
Choose any $b\in M$ with $(\varphi-1)b=y$. For $h\in N$,
$(\varphi-1)\delta b(h)=0$; also
$(\varphi-1)(z-\delta b(g_0))=0$. Thus both expressions
$\delta b(h)$ and $z-\delta b(g_0)$ lie in $j(T)$.

We construct the continuous cocycle from these values. In
$M\rtimes G'$ take the closed subgroup $C_0$ generated by
$e_0=(z,g_0)$. Every $(u',g')\in C_0$ satisfies
\[
   h u'=u'\quad(h\in N),\qquad u'-\delta b(g')\in j(T).
\]
Each condition defines a closed subgroup (use normality of $N$ in
the first and the coboundary identity in the second), and $e_0$
satisfies both. Hence $C_0$ is compact: the second condition
puts it in the continuous image of $j(T)\times G'$ by
$(w,g')\mapsto(w+\delta b(g'),g')$. The continuous map
$C_0\to\mathbb Z_p$, $(u',g')\mapsto\lambda(g')$, is onto by
compactness and density of the integers.

We show that the first component of every point over zero is zero.
Since $C_0$ is compact and procyclic in the sense of having a dense
cyclic subgroup, each $q\in C_0$ has the form $e_0^j r^p$ with
$0\le j<p$ and $r\in C_0$: the finite union of these compact sets
contains all integer powers. If $q$ is over zero then
$0=j+p\lambda(r_2)$, forcing $j=0$ and $r$ to be over zero.
Iterate, obtaining $q=r_k^{p^k}$. For such $r_k$ its second
coordinate acts trivially on its first, and its first belongs to
$j(T)$ by the two displayed conditions and $\delta b(N)\subset j(T)$.
Thus $q_1\in p^kj(T)$ for every $k$, and separation gives $q_1=0$.

Two points over the same value have the same first coordinate by
applying this to their quotient. Their common first coordinate
$d(\alpha)$ is a continuous function of $\alpha\in\mathbb Z_p$,
because the compact-to-Hausdorff surjection is a quotient map.
It satisfies $d(0)=0$, $d(1)=z$. Multiplication in $C_0$ gives
$d(\lambda(g h))=d(\lambda(g))+a_g d(\lambda(h))$:
any two choices of the group element acting on the second summand
differ by $N$, which fixes that summand. Set
\[
             j c_b(g)=d(\lambda(g))-\delta b(g).
\]
For a lift $(d(\lambda(g)),g')\in C_0$ we have
$g'^{-1}g\in N$, so $\delta b(g)-\delta b(g')\in j(T)$.
This proves integrality; the embedding of $j(T)$ gives continuity.
The displayed formula is a cocycle. It takes the required values
$-\delta b(h)$ on $N$ and $z-\delta b(g)$ on the fiber of $1$.
These two rules determine it uniquely, since the closed subgroup
generated by $N$ and $g_0$ is $G'$ (its image is compact and
contains all integers).

Replacing $b$ by $b+j(t)$ changes $c_b$ by the integral coboundary
of $-t$. It follows that $[c_b]$ is independent of $b$. Adding or
scaling pairs and their chosen lifts adds or scales their cocycles
by uniqueness. This gives the two $O$-linear maps. If
$\gamma_K=\gamma$ and the two underlying pairs agree, the full
cocycle restricted to $G_s$ satisfies the two restricted rules
for the same $b$: an identity quotient element belongs to $\Delta$,
and a lift of $\gamma$ lifts $\gamma_K$. Uniqueness gives the exact
restriction identity in integral cohomology. This includes the
defining clause of the full record under its printed dyadic
implication $p=2\Rightarrow\Delta=\{1\}$. The construction of the
maps itself allowed the full $\Delta$ in every characteristic.
\end{proof}
% END LWR5

At the per-layer choice in \texttt{ExistsLocalWachFunctor.lean},
the first three choices, the \'etale module (FET25), the Wach lattice
(WPR30 at its own premises), and the Herr maps now fit the same
original records. The remaining per-layer choice is
$\mathrm{CmpS}(\mathrm{Mo})$. It requires convergent evaluation into
$B_{\mathrm{cris}}$, the open-disc expansion and its common invariant
basis, the period map and its filtered de Rham estimates. The
calculation of level Taylor series in LWR2 does not supply these
comparison maps. The original common denominator family, integral
finite local conditions, coefficient-normalized lengths, both
orders of the signed clusters, and all shared global bounds remain
at their previously audited scope.
